% =====================================================================
%  INTRODUCTION, OBJECTIFS ET ENVIRONNEMENT DE TRAVAIL (non numéroté)
% =====================================================================
\chapter*{Introduction et objectifs}
\addcontentsline{toc}{chapter}{Introduction et objectifs}

Ce document constitue le compte rendu des travaux pratiques réalisés dans le
cadre du module «~\moduleNom~». Il présente la démarche suivie, le code
développé ainsi que les résultats obtenus et leur interprétation.

\paragraph{Objectifs du TP.}
À l'issue de cette séance, l'étudiant doit être capable de~:
\begin{itemize}[leftmargin=1.2em,itemsep=1pt]
  \item citer ici le premier objectif pédagogique du TP ;
  \item manipuler les outils et concepts abordés (langage, bibliothèque,
    algorithme\ldots) ;
  \item analyser et interpréter les résultats produits par le programme.
\end{itemize}

\paragraph{Environnement de travail.}
Indiquez ici le matériel et les logiciels utilisés~:
\begin{itemize}[leftmargin=1.2em,itemsep=1pt]
  \item \textbf{Système d'exploitation}~: Windows 11 / Linux / macOS ;
  \item \textbf{Langage / version}~: \codeinline{Python 3.12}, \codeinline{gcc 13},
    \codeinline{Java 21}\ldots;
  \item \textbf{Éditeur / IDE}~: VS Code, PyCharm, IntelliJ, Vim\ldots;
  \item \textbf{Outils complémentaires}~: Git, terminal, Jupyter Notebook.
\end{itemize}

\paragraph{Organisation du rapport.}
La \textbf{Première partie} présente le premier exercice traité. La
\textbf{Deuxième partie} détaille un second exercice. Une conclusion générale
clôt le document.

\begin{remarque}[Comment utiliser ce template]
Remplacez le texte de chaque partie par votre propre contenu. Les blocs de
code s'écrivent avec l'environnement \texttt{code} (options
\texttt{language=\ldots} et titre entre accolades), et les sorties de terminal
avec l'environnement \texttt{terminal}. Voir le fichier \texttt{README.md}
pour la liste des langages supportés et des exemples complets.
\end{remarque}
